% Figure: a bell-crank (right-angle lever) pinned at the corner, an input
% force on one arm and an output on the other, with the pin reaction. Numbers
% match the bell-crank worked example.
\begin{figure}[htbp]
\centering
\begin{tikzpicture}[
  force/.style={draw=engblue, -{Stealth}, very thick},
  rxn/.style={draw=engred, -{Stealth}, very thick},
  scale=1.0,
]
% pin at corner O
\coordinate (O) at (0,0);
\coordinate (A) at (0,2.6);   % end of vertical arm (input)
\coordinate (B) at (3.0,0);   % end of horizontal arm (output)
% the crank body: two arms meeting at O
\draw[draw=engblack, line width=2.4pt] (A) -- (O) -- (B);
% pin
\draw[draw=engblack, fill=enggray!20, thick] (O) circle (0.13);
\fill[engblack] (O) circle (0.04);
\node[font=\scriptsize, anchor=north east] at (O) {$O$};
% input force P at A, horizontal (pushing right)
\draw[force] (A) -- ++(1.7,0) node[anchor=west, font=\small, engblue] {$P$};
\node[font=\scriptsize, anchor=east] at (A) {$A$};
% arm lengths
\draw[draw=enggray, thin, <->] (-0.35,0) -- (-0.35,2.6);
\node[font=\scriptsize, enggray, anchor=east] at (-0.38,1.3) {$a$};
\draw[draw=enggray, thin, <->] (0,-0.35) -- (3.0,-0.35);
\node[font=\scriptsize, enggray, anchor=north] at (1.5,-0.4) {$b$};
% output force Q at B, vertical (up)
\draw[force] (B) -- ++(0,1.7) node[anchor=south, font=\small, engblue] {$Q$};
\node[font=\scriptsize, anchor=north west] at (B) {$B$};
% pin reaction (schematic, at O, up-left)
\draw[rxn] (O) -- ++(-1.2,1.0) node[anchor=south east, font=\small, engred] {$R_O$};
\end{tikzpicture}
\caption[A bell-crank lever]{The chosen body is the bell-crank, a rigid
right-angle lever pinned at the corner $O$. An input push $P$ acts
horizontally at the end $A$ of the arm of length $a$; the output $Q$ leaves
vertically at the end $B$ of the arm of length $b$. The pin supplies two
reaction components, drawn here as a single resultant $R_O$. Summing moments
about $O$ kills the pin reaction, because it passes through $O$, and leaves
one equation relating input to output: $P\,a = Q\,b$. The crank turns the
line of a force through a right angle and scales its magnitude by the arm
ratio $a/b$, which is why a bell-crank is the fitting of choice wherever a
control force must change direction, from a bicycle brake to an aircraft
control run. The pin then carries the vector sum of the two arm forces, a
load neither arm alone reveals.}
\label{fig:vol03ch01:bell-crank}
\end{figure}
